\documentstyle[% 


righttag% 


,11pt% 


,amssymb% 


,russian%               remove in Eng. 


,izvru%                 Set izven% in Eng. 


% ,izvru-ks             for brief communications 


]{amsart} 


 


% 


 


\newcommand{\M}{\mathrm} 


\newcommand{\tts}[1]{} 


 


\theoremstyle{plain} 


\theorembodyfont{\it} 


\newtheorem{thmnonum}{\hskip\parindent Теорема} 


\newtheorem{lemnonum}{\hskip\parindent Лемма} 


\renewcommand{\thethmnonum}{} 


\renewcommand{\thelemnonum}{} 


 


 


%\hoffset=-15mm%                  For Eng. 


\setcounter{page}{52} 


\god{2002} 


\nomer{11~(486)} %                        Remove in Eng. 


%\nomer{Vol.\,46, No.\,11}%               Set in Eng. 


%\str{\pageref{begin}--\pageref{end}}%  Set in Eng. 


\udk{514.4} 


 


\author{\it М.А.\,Микенберг} 


% NB:   ^^ remove in Eng. 


\title{Некоторые топологические свойства  


многообразий над локальной алгеброй,  


допускающих голоморфные вложения} 


 


\date{} 


%\pagestyle{myheadings}%         Set in Eng. 


\pagestyle{plain} %              Remove in Eng. 


\begin{document} 


%\thispagestyle{plain}%          Set in Eng. 


\maketitle 


\label{begin} 


 


 


 


\section{Введение} 


 


Обозначим через $\cal A$ локальную алгебру над полем вещественных 


чисел $\M R$. Всюду далее локальная алгебра $\cal A$ будет конечномерной, 


ассоциативной, коммутативной, унитальной алгеброй. Обозначим через 


$\M I$ радикал алгебры $\cal A$. Как векторное пространство 


над полем $\M R$ алгебра $\cal A$ изоморфна прямой сумме $\M R\oplus\M 


I$. Обозначим 


через $\M n$ размерность алгебры $\cal A$ над полем $\M R$, $\cal 


A\approx\M R^{\M n}$. Обозначим 


через $\M M$ конечномерное многообразие над $\cal A$. Напомним 


(см. [1], [2]), что многообразие $\M M$ называется многообразием над 


алгеброй $\cal A$, если у него существует атлас $\{(\M U_{\M 


i},\varphi_{\M i})\}_{\M i\in\M I}$ такой, что 


$\varphi(\M U_{\M i})\subset\cal A^{\M m}$ (где $\M m$ одно и то же для 


всех $\M i$), а функции перехода 


$\varphi_{\M j}\circ\varphi^{-1}_{\M i}$ \ $\cal A$-голоморфны. Атлас 


$\{(\M U_{\M i},\varphi_{\M i})\}_{\M i\in\M I}$ будем называть 


$\cal A$-атласом. (Напомним, что отображение $f:(\M U\subset\cal A^{\M 


p})\to\cal A^{\M q}$, 


$\M U$ открыто в $\cal A^{\M p}$, называется $\cal A$-голоморфным, если 


оно принадлежит 


классу $C^\infty$ как отображение над $\M R$, если отождествить $\cal A$ с 


$\M R^{\M n}$ ($\M n$ --- размерность алгебры $\cal A$ над полем $\M R$) и 


производная $f'$ \ 


$\cal A$-линейна.) Число $\M m$ называется размерностью многообразия 


$\M m$ 


над $\cal A$. Любое многообразие над алгеброй $\cal A$ можно рассматривать 


как вещественное многообразие. Размерность многообразия $\M M$ над 


$\M R$ равна $\M {mn}$. В данной статье изучаются голоморфные вложения 


многообразия $\M M$ над локальной алгеброй $\cal A$ в некоторую декартову 


степень $\cal A^{\M r}$ алгебры $\cal A$. 


 


\section{Некоторые топологические свойства многообразий 


над локальной алгеброй $\cal A$, голоморфно вложенных в некоторую 


декартову 


степень $\cal A^{\M r}$ алгебры $\cal A$} 


 


Пусть далее $\M M$ --- многообразие размерности $\M m$ над локальной 


алгеброй $\cal A$, которая имеет размерность $\M n$ над $\M R$. 


Отображение 


$\M F:\M M\to\cal A^{\M r}$ называется $\cal A$-голоморфным вложением 


многообразия $\M M$ 


над локальной алгеброй $\cal A$ в $\cal A^{\M r}$, если $\M F$ задает 


вложение 


класса $C^\propto$ многообразия $\M M$, рассматриваемого как вещественное 


многообразие, в $\M R^{\M {rn}}$ ($\cal A^{\M r}\approx\M R^{\M {rn}}$) и 


отображение $\M F$ \ 


$\cal A$-голоморфно. Последнее означает, что многообразие $\M M$ обладает 


$\cal A$-атласом $\{(\M U_{\M i},\varphi_{\M i})\}_{\M {i\in I}}$ таким, 


что отображения 


$\M F\circ\varphi^{-1}_{\M i}:(\varphi_{\M i}(\M U_{\M i})\subset\cal 


A^{\M m})\to\cal A^{\M r}$ \ $\cal A$-голоморфны. Рассмотрим нижний 


центральный ряд радикала $\M I$ алгебры $\cal A$ 


$$ 


\M I\supset\M I^2\supset\M I^3\supset\dotsb\supset 


\M I^{\M q-1}\supset\M I^{\M q}\supset\M I^{\M{q+1}}=0. 


$$ 


 


Рассмотрим идеал $\M I^{[\M q/2]+1}$. Здесь через $[\M q/2]$ обозначена 


целая часть числа $\M q/2$. Пусть размерность идеала $\M I^{[\M q/2]+1}$ 


над полем $\M R$ равна $\M k$. Обозначим через $\M I^{[\M 


q/2]+1}\cdot\M T_{\M x}\M M$ 


подпространство в касательном пространстве $\M T_{\M x}\M M$ многообразия 


$\M M$ в точке $\M x$, состоящее из векторов вида $z=\sum\limits^{\M 


m}_{i=1}\alpha_i z_i$, 


где $\alpha_i\in\M I^{[\M q/2]+1}$, $z_i\in\M T_{\M x}\M M$. Размерность 


подпространства 


$\M I^{[\M q/2]+1}\cdot\M T_{\M x}\M M$ над полем $\M R$ равна $\M{mk}$. 


 


\begin{lemnonum} 


Для любой голоморфной системы координат $(U,\varphi)$ на $\M M$ и для 


любых $u,v\in(\M I^{[\M q/2]+1})^{\M m}\subset\M A^{\M m}$ 


справедливо $(\M F\circ\varphi^{-1})''(u,v)=0$. 


\end{lemnonum} 


 


\begin{pf} 


Рассмотрим голоморфную систему координат $(\M U,\varphi)$ 


в окрестности точки $\M x$. Она задает изоморфизм между $\M T_{\M x}\M M$ и 


$\cal A^{\M m}\approx\M R^{\M {mn}}$ и 


между подпространством $\M I^{[\M q/2]+1}\cdot\M T_{\M x}\M M$ и 


декартовой степенью $(\M I^{[\M q/2]+1})^{\M m}$ 


идеала $\M I^{[\M q/2]+1}$. Пусть $\wt{\M F}=\M F\circ\varphi^{-1}$, $\wt{\M 


x}=\varphi(\M x)$. Вторая производная $\wt{\M F}{}''(x)$ отображения 


$\wt{\M F}$ в точке $\wt{\M x}$ тождественно равна нулю на 


декартовом произведении $(\M I^{[\M q/2]+1})^{\M m}$ идеала $\M I^{[\M 


q/2]+1}$. Действительно, 


пусть $u,v\in(\M I^{[\M q/2]+1})^{\M m}$, тогда 


$u=\sum\limits^{\M m}_{i=1}\alpha_i u_i$, $v=\sum\limits^{\M 


m}_{i=1}\beta_i v_i$, где $u_i,v_i\in\cal A^{\M m}\approx\M R^{\M {mn}}$, 


$\alpha_i,\beta_i\in\M I^{[\M q/2]+1}$. 


Так как отображение $\M F$ \ $\cal A$-голоморфно,


то $\wt{\M F}{}''(\wt{\M x})$ --- симметрическая $\cal A$-билинейная форма.


Таким образом, 


$$ 


\wt{\M F}{}''(\wt{\M x})(u,v)=\wt{\M F}{}''(\wt{\M x})\bigg(\, 


\sum^m_{i=1}\alpha_i u_i,\


\sum^{\M m}_{i=1}\beta_i v_i\bigg)=\sum^{\M m}_{i=1} 


\sum^{\M m}_{j=1}\alpha_i\beta_j 


\wt{\M F}{}''(\wt{\M x})(u_i,v_j). 


$$ 


(Здесь через $\wt {\M F}{}''(\wt{\M x})(u,v)$ обозначено значение второй 


производной 


отображения $\wt{\M F}$ в точке $\wt{\M x}$ на элементах $u$ и $v$.) Но 


$\alpha_i\beta_j=0$, а значит, $\wt{\M F}{}''(\wt{\M x})(u,v)=0$. 


\end{pf} 


 


Пусть $\M F:\M M\to\cal A^{\M r}$ --- $\cal A$-голоморфное вложение 


многообразия $\M M$. 


Введем в $\cal A^{\M r}\approx\M R^{\M{rn}}$ евклидову структуру. 


Скалярное 


произведение в $\cal A^{\M r}\approx\M R^{\M {rn}}$ будем обозначать через 


$(a,b)$, 


где $a,b\in\cal A^{\M r}\approx\M R^{\M {rn}}$. Для почти всех векторов 


(т.\,е. для всех 


векторов, кроме множества меры нуль (см. [3], теорема 6.6)) 


$\M p\in\cal A^{\M r}\approx\M R^{\M {rn}}$ функция $\M L_{\M p}:\M 


M\to\M R$, где 


$$ 


\M L_{\M p}(\M x)=(\M F(\M x)-\M p,\M F(\M x)-\M p), 


$$ 


имеет невырожденные критические точки. Имеем 


$$ 


\M L'_{\M p}(\M x)=2(\M F'(\M x),\M F(\M x)-\M p), 


$$ 


где $\M L'_{\M p}(\M x)$ --- первая производная отображения $\M L_{\M 


p}$ в точке $\M x$. 


Пусть $\M x_0$ --- критическая точка отображения $\M L_{\M p}$. Введем 


голоморфную систему координат в окрестности точки $\M x_0$, 


в которой будем проводить дальнейшие вычисления. Имеем 


$$ 


\M L'_{\M p}(\M x_0)=2(\M F'(\M x_0), 


\M F(\M x_0)-\M p)=0 


$$ 


и 


$$ 


\M L''_{\M p}(\M x_0)=2((\M F'(\M x_0),\M F'(\M x_0))+ 


(\M F''(\M x_0),\M F(\M x_0)-\M p)), 


$$ 


где $\M L''_{\M p}(\M x_0)$ --- вторая производная отображения $\M 


L_{\M p}$ в критической 


точке $\M x_0$. \ $\M L'_{\M p}(\M x)$ --- $\M R$-линейная форма, 


действующая в $\M T_{\M x}\M M$, а если 


$\M x_0$ --- критическая точка отображения $\M L_{\M p}$, то 


$\M L''_{\M p}(\M x_0)$ --- $\M R$-билинейная симметрическая форма, 


действующая в 


касательном пространстве $\M T_{\M x_0}\M M$. Пусть $(\M F'(\M x),\M 


F'(\M x))$ --- первая квадратичная 


форма многообразия $\M M$ (вложенного в $\M R^{\M {rn}}\approx\cal 


A^{\M r}$). Если $\M x_0$ --- 


критическая точка отображения $\M L_{\M p}$, то $(\M F''(\M x_0),\M 


F(\M x_0)-\M p)$ --- 


вторая квадратичная форма многообразия $\M M$ в точке $\M x_0$, 


вычисленная в направлении вектора $\M F(\M x_0)-\M p$. Если $\M x_0$ --- 


невырожденная 


критическая точка функции $\M L_{\M p}$, то $\M L'_{\M p}(\M x_0)=0$, а 


билинейная 


форма $\M L''_{\M p}(\M x_0)$ невырождена (см. [3], [4]). 


 


 


\begin{thmnonum} 


Если $\M F:\M M\to\cal A^{\M r}$ --- $\cal A$-голоморфное вложение 


многообразия $\M M$ 


над локальной алгеброй $\cal A$ такое, что образ $\M F(\M M)$ 


замкнут в $\cal A^{\M r}\approx\M R^{\M {rn}}$, то $\M M$ имеет 


гомотопический тип 


клеточного комплекса размерности, не больше чем $\M m(\M n-\M k)$, где 


$\M k$ --- 


размерность идеала $\M I^{[\M q/2]+1}$ над $\M R$. 


\end{thmnonum} 


 


\begin{pf} 


Выберем точку $\M p\in\cal A^{\M r}$ так, что функция 


$$ 


\M L_{\M p}(\M x)=(\M F(\M x)-\M p, 


\M F(\M x)-\M p) 


$$ 


имеет невырожденные критические точки. Тогда $\M M$ имеет гомотопический 


тип клеточного комплекса, поскольку компактны множества 


$\M M^a=\M L_{\M p}^{-1}([0,a])$ (см. [3], теорема 3.5; а также [4]). 


При этом каждой критической точке $\M x_0$ с индексом $\lambda$ функции 


$\M L_{\M p}$ соответствует одна клетка размерности $\lambda$. В 


критической 


точке $\M x_0$ в силу леммы билинейная форма 


$$ 


\M L''_{\M p}(\M x_0)=2((\M F'(\M x_0),\M F'(\M x_0))+ 


(\M F''(\M x_0),\M F(\M x_0)-\M p)) 


$$ 


положительно определена на подпространстве пространства $\M T_{\M 


x_0}\M M$ 


размерности $\M {mk}$. Значит, билинейная форма $\M L''_{\M p}(\M x_0)$ 


отрицательно определена 


на подпространстве пространства $\M T_{\M x_0}\M M$ размерности, не выше 


чем $\M m(\M n-\M k)$. 


\end{pf} 


 


\begin{cor} 


В условиях теоремы $\M H_{\M i}(\M M,\M Z)=0$, если 


$\M i>\M m(\M n-\M k)$. (Здесь через $\M H_{\M i}(\M M,\M Z)$ 


обозначена $\M i$-я группа 


целочисленных гомологии многообразия $\M M$.) 


\end{cor} 


 


\begin{cor} 


Компактное многообразие над локальной алгеброй 


$\cal A$ не допускает $\cal A$-го\-ло\-морф\-ное вложение в $\cal A^{\M r}$


ни для какого $\M r$. 


\end{cor} 


 


\begin{cor} 


Если в условиях теоремы $\cal A$ --- алгебра дуальных 


чисел, то многообразие $\M M$ имеет гомотопический тип клеточного 


комплекса размерности, не больше чем $\M m$. Следовательно, 


$\M H_{\M i}(\M M,\M Z)=0$, если $\M i>\M m$. 


\end{cor} 


 


 


\begin{thebibliography}{9} 


 


\bibitem{1} 


Вишневский В.В., Широков А.П., Шурыгин В.В. {\em Пространства 


над алгебрами}. -- Казань: Изд-во Казанск. ун-та, 1985. -- 264 c.


\bibitem{2} 


Шурыгин В.В. {\em Гладкие многообразия над локальными алгебрами, 


их применение в дифференциальной геометрии высшего 


порядка}. Дис.~$\dotsc$ докт. физ.-матем. наук. -- 


Казань, 1998. -- 128~c.


\bibitem{3} 


Милнор Дж. {\em Теория Морса}. -- М.: Мир, 1965.


\bibitem{4} 


Постников М.М. {\em Введение в теорию Морса}. -- М.: Мир, 


1965.


 


\end{thebibliography} 


 


\vskip 2cm 


 


\post{\it Самарский государственный}{\it Поступила} 


\post{\it педагогический университет}{17.05.2001} 


 


\label{end} 


\end{document} 


 


